\documentclass[11pt,onecolumn]{article}
\usepackage[T1]{fontenc}

\usepackage{amsmath,amsfonts,amssymb}
\usepackage[margin=1in]{geometry}
\usepackage{hyperref}
\usepackage{parskip}
\usepackage{titlesec}

\titleformat{\section}{\large\bfseries}{}{0em}{}[\titlerule]
\titlespacing{\section}{0pt}{12pt}{6pt}

\title{\textbf{Implementation Notes}\\
\large Speech Understanding --- Programming Assignment 2\\[4pt]
\normalsize One Non-Obvious Design Choice per Task}
\author{Roll No: M25CSA028 \quad April 2026\\
\small \url{https://github.com/kingkenche/Speech-Understanding/tree/Assignment-2}}
\date{}

\begin{document}

\maketitle

\section{Task 1.1 --- Multi-Head LID: Frame-Level CNN over Utterance-Level Classification}

\textbf{Design Choice:} The LID model (\texttt{FrameLIDNet}) operates on 800\,ms sliding windows
(40\,$\times$\,20\,ms frames) with a stride of one frame, rather than classifying whole utterances.

\textbf{Rationale:} Code-switching in Hinglish happens sub-utterance --- often mid-sentence
(e.g., ``yeh \textit{stochastic} process hai''). Utterance-level classifiers average out these
transitions, producing a single language label that is wrong for both segments. By sliding a CNN
over log-mel frames with a receptive field of $\sim$800\,ms, the model produces a per-frame
language probability sequence. Switch boundaries are detected as the first frame where the
predicted class changes, achieving the $\pm$200\,ms precision required by the assignment. A
multi-head variant applies two separate linear classification heads to the shared encoder --- one
per language --- enabling independent calibration of English and Hindi posteriors before the argmax
decision.

\section{Task 1.2 --- Constrained Decoding: Additive Logit Bias over Hard Vocabulary Constraints}

\textbf{Design Choice:} The \texttt{NGramLogitBiasProcessor} adds a \textit{scaled n-gram
probability} to each candidate token's logit rather than zeroing out non-technical tokens (hard
constraint / vocabulary whitelisting).

\textbf{Rationale:} Hard vocabulary constraints force Whisper to pick technical terms even when
acoustic evidence strongly suggests a common word, causing hallucinations. Additive biasing
\[
\tilde{z}_t^{(v)} = z_t^{(v)} + \alpha \cdot \log P_{\text{ngram}}(x_v \mid h)
\]
instead \textit{probabilistically steers} the model: if the n-gram model assigns high probability
to ``cepstrum'' given the context ``mel frequency'', the bias nudges beam-search ranking upward
without overriding acoustic evidence. The scale factor $\alpha = 4.0$ was chosen so that a
trigram probability of 0.1 contributes approximately the log-space equivalent of one acoustic
frame of evidence --- a sweet spot that boosts recall of technical terms without degrading
fluency.

\section{Task 2.1 --- Hinglish G2P: Greedy Digraph Matching over Neural Lookup Tables}

\textbf{Design Choice:} The \texttt{\_word\_to\_ipa()} function uses greedy left-to-right
\textit{digraph matching} rather than a trained G2P neural model for the Devanagari-romanized
portion of Hinglish.

\textbf{Rationale:} No publicly available G2P model handles code-switched Romanized Hindi
reliably. Hindi romanization is not standardized --- ``dh'', ``bh'', ``ph'' are aspirated stops
mapping to different IPA symbols than their component characters. A greedy digraph scan (check
2-character pair first, fall back to single character) reliably handles all standard ITRANS-style
romanizations without any training data. For English tokens, standard grapheme-to-IPA mapping
handles the ASCII portion. The combined system correctly converts
``bhasha'' $\rightarrow$ /bh-aa-sh-a/ and ``deep'' $\rightarrow$ /d-ii-p/
without any trained model, making it fully reproducible with zero labelled data.

\section{Task 3.2 --- Prosody Warping: Energy-Scaled DTW over Vocoder Pitch Shifting}

\textbf{Design Choice:} Instead of vocoder-based pitch shifting (PSOLA), \texttt{apply\_prosody\_warping()} applies DTW-aligned $F_0$ and energy
ratios directly to the waveform amplitude envelope.

\textbf{Rationale:} Vocoder-based pitch shifting requires resynthesis from acoustic parameters,
introducing a full vocoder round-trip that degrades perceptual quality. The surrogate approach
--- warping the energy envelope via DTW alignment indices and applying the $F_0$ ratio as a mild
amplitude modulation --- transfers the \textit{teaching rhythm} (louder emphasis on key words,
pauses between concepts) without altering the fundamental perceptual character of the synthesized
voice. For TTS output, the energy envelope encodes prosodic emphasis more reliably than raw $F_0$
because synthesized $F_0$ is already smooth; modulating it again introduces artifacts. This avoids
the MOS degradation associated with double-synthesis while still preserving the temporal pattern
of the professor's lecture delivery.

\section{Task 4.1 --- Anti-Spoofing: Shallow TDNN over SVM on LFCC Features}

\textbf{Design Choice:} The \texttt{LFCCTDNN} model uses two 1D convolutional layers (equivalent
to a shallow Time Delay Neural Network) with adaptive average pooling rather than a linear SVM.

\textbf{Rationale:} LFCC feature sequences vary in length across utterances; an SVM requires
fixed-size input. Padding to maximum length and flattening loses temporal structure and creates a
high-dimensional sparse input where most classifiers fail. The TDNN architecture
(\texttt{Conv1D} $\rightarrow$ \texttt{ReLU} $\rightarrow$ \texttt{Conv1D} $\rightarrow$
\texttt{AdaptiveAvgPool1d(1)}) naturally handles variable-length LFCC matrices by collapsing the
time axis to a single 96-dimensional vector via global average pooling, which a linear head then
classifies. This is simpler than attention and achieves comparable EER to SVM on the ASVspoof
benchmark for the binary bona-fide vs.\ spoof case, while remaining fully differentiable for
FGSM gradient-based adversarial analysis in Task~4.2.

\end{document}
